% ---------------------------------------------------------------------------------------

[!] UOC - CSD - HY110 - #1 - csd5748

% ---------------------------------------------------------------------------------------

\documentclass{article}
\usepackage[utf8]{inputenc}  % Για υποστήριξη της UTF-8 κωδικοποίησης.
\usepackage[greek]{babel}    % Για να γράφουμε στα ελληνικά.
\usepackage{amsmath}         % Για την εντολή \xRightarrow{}.
\usepackage{amsfonts}        % For \mathbb.
\usepackage{amssymb} 
\usepackage[a4paper, left=1cm, right=1cm, top=2cm, bottom=2cm]{geometry}

\begin{document}

% ---------------------------------------------------------------------------------------

\newpage
\noindent \textbf{ }

% ---------------------------------------------------------------------------------------

\newpage

\vspace{3cm}

\begin{center}
\noindent \boxed{\textbf{· 1\textsuperscript{η} Άσκηση:} \textbf{Υπολογισμός Ορίων !}}
\end{center}

\vspace{1.25cm}

\noindent \textbf{· Άσκηση 1\textsuperscript{η} - Ερώτημα 1\textsuperscript{ο}:} \hfill \\[1ex]

\[
    \lim_{x \to 0^+} \left( \frac{\sqrt{5} - \sqrt{4x^2 + 9x + 5}}{x} \right) \hspace{0.125cm} = \hspace{0.175cm} \textlatin{?} 
\]

\vspace{1cm}

\noindent \textbf{· Η Λύση μου:} \hfill

\begin{align*}
    L_1 \hspace{0.125cm} & = \hspace{0.1cm} \lim_{x \to 0^+} \left( \frac{\sqrt{5} - \sqrt{4x^2 + 9x + 5}}{x} \right) = \hspace{0.125cm} \frac{0}{0} \hspace{0.175cm} \text{“A.M.”} \\[5ex]
    & \mathbf{\Rightarrow} \hspace{0.15cm} \lim_{x \to 0^+} \left( \frac{(\sqrt{5} - \sqrt{4x^2 + 9x + 5}) \cdot (\sqrt{5} + \sqrt{4x^2 + 9x + 5})}{x \cdot (\sqrt{5} + \sqrt{4x^2 + 9x + 5})} \right) \\[5ex]
    & = \hspace{0.1cm} \lim_{x \to 0^+} \left( \frac{(\sqrt{5})^2 - \left(\sqrt{4x^2 + 9x + 5}\right)^2}{x \cdot (\sqrt{5} + \sqrt{4x^2 + 9x + 5})} \right) \\[5ex]
    & = \hspace{0.1cm} \lim_{x \to 0^+} \left( \frac{-4x^2 - 9x \pm 0}{x \cdot (\sqrt{5} + \sqrt{4x^2 + 9x + 5})} \right) \\[5ex]
    & = \hspace{0.1cm} \lim_{x \to 0^+} \left( \frac{x \cdot (-4x - 9)}{x \cdot (\sqrt{5} + \sqrt{4x^2 + 9x + 5})} \right) \\[5ex]
    & = \hspace{0.1cm} \lim_{x \to 0^+} \left( \frac{-4x - 9}{\sqrt{5} + \sqrt{4x^2 + 9x + 5}} \right) \\[5ex]
    & = \hspace{0.1cm} \frac{-4(0^+) - 9}{\sqrt{5} + \sqrt{4 (0^+)^2 + 9 (0^+) + 5}} \\[5ex]
    & = \hspace{0.1cm} \frac{-9}{\sqrt{5} + \sqrt{5}} \\[5ex]
    & \hspace{0.125cm} = \hspace{0.125cm} - \frac{9}{2\sqrt{5}} \\[5ex]
    & \hspace{0.125cm} = \hspace{0.125cm} \boxed{-\frac{9\sqrt{5}}{10}}
\end{align*}

% ---------------------------------------------------------------------------------------

\newpage

\noindent \textbf{· Άσκηση 1\textsuperscript{η} - Ερώτημα 2\textsuperscript{ο}:} \hfill

\[
    \lim_{x \to \infty} \left( \frac{x + \cos{(x)}}{2x + 6 - 3\sin{(x)}} \right) \hspace{0.125cm} = \hspace{0.175cm} \textlatin{?}
\]

\vspace{0.5cm}

\noindent \textbf{· Η Λύση μου:} \hfill

\[
    L_2 \hspace{0.125cm} = \hspace{0.1cm} \lim_{x \to \infty} \left( \frac{x + \cos{(x)}}{2x + 6 - 3\sin{(x)}} \right) = \hspace{0.125cm} \frac{\infty}{\infty} \hspace{0.175cm} \text{“A.M.”} \\[5ex]
\]
    
\[
    \mathbf{\Rightarrow} \hspace{0.15cm} 
    \lim_{x \to \infty} \left( 
    \frac{x( 1 + \frac{\cos{(x)}}{x})}{x( 2 + \frac{6}{x} - 3\frac{\sin(x)}{x})} \right) 
    \hspace{0.15cm} = \hspace{0.15cm} 
    \lim_{x \to \infty} \left( 
    \frac{1 + \frac{\cos{(x)}}{x}}{2 + \frac{6}{x} - 3\frac{\sin(x)}{x}} \right) 
    \hspace{0.15cm} = \hspace{0.15cm} 
    \lim_{x \to \infty} \left( 
    \frac{1 + 0}{2 + 0 - 3 \cdot 0} \right) 
    \hspace{0.15cm} = \hspace{0.15cm} 
    \boxed{\frac{1}{2}} 
    \\[5ex]
\]

\[
    \textbf{αφού:} \hspace{0.25cm} \lim_{x \to \infty} \left( 
    \frac{6}{x} \right) 
    \hspace{0.15cm} = \hspace{0.15cm} 
    \frac{6}{\infty} \hspace{0.15cm} = \hspace{0.15cm} 0 \hspace{9cm} { } 
    \\[5ex]
\]

\[
    \textbf{και:} \hspace{0.25cm} 
    \lim_{x \to \infty} \left( \frac{\cos{x}}{x} \right) 
    \hspace{0.15cm} \textbf{:} \hspace{0.25cm} \hspace{10.5cm} { }  
    \\[1.5ex]
\]

\[
    \textbf{είνα: } \hspace{0.35cm}
    -1 \leq \cos(x) \leq 1 
    \hspace{0.5cm} 
    \overset{x \neq 0}{\Rightarrow}
    \hspace{0.5cm} 
    \frac{-1}{x} \leq \frac{\cos(x)}{x} \leq \frac{1}{x} 
    \hspace{5.55cm} { }
    \\[1.5ex]
\]

\[
    \overset{\textbf{ΘΠ.}}{\Rightarrow}
    \hspace{0.5cm} 
    \lim_{x \to \infty} \left( -\frac{1}{x} \right) 
    \hspace{0.15cm} = \hspace{0.15cm} 
    \lim_{x \to \infty} \left( \frac{1}{x} \right) 
    \hspace{0.15cm} = \hspace{0.15cm} 
    \lim_{x \to \infty} \left( \frac{\cos(x)}{x} \right) 
    \hspace{0.15cm} = \hspace{0.15cm} 0 
    \hspace{4.5cm} { }
    \\[3.5ex]
\]

\[
    \textbf{και ομοίως για το:} \hspace{0.25cm} 
    \lim_{x \to \infty} \left( \frac{\sin{x}}{x} \right) 
    \hspace{0.15cm} = \hspace{0.25cm} 0 \hspace{7cm} { }  
    \\[27.5ex]
\]

% ---------------------------------------------------------------------------------------

\noindent \textbf{· Άσκηση 1\textsuperscript{η} - Ερώτημα 3\textsuperscript{ο}:} \hfill

\[
    \lim_{x \to 5} \left( \frac{5 - x}{6 - \sqrt{x^2 + 10}} \right) 
    \hspace{0.125cm} = \hspace{0.175cm} \textlatin{?}
\]

\vspace{0.5cm}

\noindent \textbf{· Η Λύση μου:} \hfill

\[
    L_3 \hspace{0.125cm} = \hspace{0.1cm} 
    \lim_{x \to 5} \left( \frac{5 - x}{6 - \sqrt{x^2 + 10}} \right) = 
    \hspace{0.125cm} \frac{0}{6 - \sqrt{35}} 
    \hspace{0.125cm} = \hspace{0.1cm} \boxed{0}
    \\[0.5ex]
\]

\newpage

% ---------------------------------------------------------------------------------------

\noindent \textbf{· Άσκηση 1\textsuperscript{η} - Ερώτημα 4\textsuperscript{ο}:} \hfill

\[
    \lim_{x \to \infty} \left( \frac{3x^2 + 5}{6x^2 + 9} \right) 
    \hspace{0.125cm} = \hspace{0.175cm} \textlatin{?}
\]

\vspace{0.5cm}

\noindent \textbf{· Η Λύση μου:} \hfill

\[
    L_4 \hspace{0.125cm} = \hspace{0.1cm} 
    \lim_{x \to \infty} \left( \frac{3x^2 + 5}{6x^2 + 9} \right) 
    \hspace{0.125cm} \overset{\textbf{Ιδ.}}{\Rightarrow} \hspace{0.125cm}
    \hspace{0.125cm} \lim_{x \to \infty} \left( \frac{3x^2}{6x^2} \right) =
    \hspace{0.125cm} \lim_{x \to \infty} \left( \frac{3}{6} \right) =
    \hspace{0.125cm} \frac{3}{6}
    \hspace{0.125cm} = \hspace{0.125cm} \boxed{\frac{1}{2}}
    \\[26ex]
\]

% ---------------------------------------------------------------------------------------

\noindent \textbf{· Άσκηση 1\textsuperscript{η} - Ερώτημα 5\textsuperscript{ο}:} \hfill

\[
    \lim_{x \to 0} \left( \frac{\sin(6x)}{\sin(3x)} \right) 
    \hspace{0.125cm} = \hspace{0.175cm} \textlatin{?}
\]

\vspace{0.5cm}

\noindent \textbf{· Η Λύση μου:} \hfill

\[
    L_5 \hspace{0.125cm} = \hspace{0.1cm} 
    \lim_{x \to 0} \left( \frac{\sin(6x)}{\sin(3x)} \right)  
    \hspace{0.125cm} \overset{x \neq 0}{\Rightarrow} \hspace{0.125cm}
    \hspace{0.125cm} \lim_{x \to 0} \left( \frac{\frac{\sin(6x)}{x}}{\frac{\sin(3x)}{x}} \right)
    \hspace{0.125cm} \overset{\textbf{Ιδ.}}{\Rightarrow} \hspace{0.125cm}
    \hspace{0.155cm} \frac{\lim_{x \to 0} \left( \frac{\sin(6x)}{x} \right)}
    {\lim_{x \to 0} \left( \frac{\sin(3x)}{x} \right)}  =
    \hspace{0.175cm} \frac{\textlatin{}{I1}}{\textlatin{}{I2}}
    \hspace{0.175cm} \Rightarrow
    \\[2.5ex]
\]
    
\[
    \textlatin{I1:} \hspace{0.255cm}
    \lim_{x \to 0} \left( \frac{\sin(6x)}{x} \right) =
    \hspace{0.125cm} \lim_{x \to 0} \left( \frac{\sin(6x)}{x} \textbf{ · } \frac{6}{6} \right) =
    \hspace{0.125cm} \lim_{x \to 0} \left( 6 \textbf{ · } \frac{\sin(6x)}{6x}\right) 
    \hspace{0.125cm} \Rightarrow 
    \hspace{1.175cm} { }
    \\[0.55cm]
\]

\[
    \textbf{Θέτω: } \hspace{0.125cm} \textbf{ω} = 6x \hspace{0.350cm} \textbf{και έχουμε:} 
    \hspace{0.125cm} \lim_{x \to 0} \left( 6 \textbf{ · } \frac{\sin(\textbf{ω})}{\textbf{ω}}\right) 
    \hspace{0.125cm} \Rightarrow
    \hspace{3.175cm} { }
    \\[0.55cm]
\]

\[
    \textbf{Όμως, ξέρουμε: } \hspace{0.125cm} 
    \lim_{x \to 0} \left( \frac{\sin(x)}{x} \right) = 1
    \hspace{0.325cm} \textbf{'αρα:} \hspace{0.325cm} 
    \Rightarrow \hspace{0.255cm} 
    \textlatin{I1 } = \hspace{0.125cm} 6 \textbf{ · } 1 \hspace{0.125cm} = \hspace{0.125cm} 6
    \hspace{0.85cm} { }
    \\[0.75cm]
\]

\[
    \textlatin{I2:} \hspace{0.255cm}
    \lim_{x \to 0} \left( \frac{\sin(3x)}{x} \right) =
    \hspace{0.125cm} \lim_{x \to 0} \left( \frac{\sin(3x)}{x} \textbf{ · } \frac{3}{3} \right) =
    \hspace{0.125cm} \lim_{x \to 0} \left( 3 \textbf{ · } \frac{\sin(3x)}{3x}\right) 
    \hspace{0.125cm} \Rightarrow
    \hspace{1.05cm} { }
    \\[0.55cm]
\]

\[
    \textbf{Και, ομοίως είναι: } \hspace{0.125cm} 
    \Rightarrow \hspace{0.255cm} 
    \textlatin{I2 } = \hspace{0.125cm} 3 \textbf{ · } 1 \hspace{0.125cm} = \hspace{0.125cm} 3
    \hspace{4.575cm} { }
    \\[0.75cm]
\]

\[
    \textbf{Οπότε: } \hspace{0.125cm} 
    \Rightarrow \hspace{0.255cm} 
    L_5 \hspace{0.125cm} = \hspace{0.125cm} \frac{6}{3}
    \hspace{0.125cm} = \hspace{0.125cm} \boxed{2}
    \hspace{6.85cm} { }
    \\[2.5cm]
\]

\newpage

% ---------------------------------------------------------------------------------------

\vspace{2.5cm}

\noindent \textbf{· Άσκηση 1\textsuperscript{η} - Ερώτημα 6\textsuperscript{ο}:} \hfill

\vspace{1.5cm}

\[
    \lim_{x \to \infty} \left( \frac{x + \cos(x) + 3\sqrt{x}}{x + \cos(x)} \right) 
    \hspace{0.125cm} = \hspace{0.175cm} \textlatin{?}
\]

\vspace{1.5cm}

\noindent \textbf{· Η Λύση μου:} \hfill

\vspace{1.5cm}

\[
    L_6 \hspace{0.125cm} = \hspace{0.1cm} 
    \lim_{x \to \infty} \left( \frac{x + \cos(x) + 3\sqrt{x}}{x + \cos(x)} \right)  
    \hspace{0.125cm} \overset{x \neq 0}{=} \hspace{0.125cm}
    \lim_{x \to \infty} \left( \frac{x \textbf{ · } (1 + \frac{\cos(X)}{x} + \frac{3\sqrt{x}}{x})}
    {x \textbf{ · } (1 + \frac{\cos(X)}{x})} \right)
    \hspace{0.125cm} = \hspace{0.125cm}
    \lim_{x \to \infty} \left( \frac{1 + \frac{\cos(X)}{x} + \frac{3\sqrt{x}}{x}}
    {1 + \frac{\cos(X)}{x}} \right)
    \hspace{0.125cm} \Rightarrow
    \\[0.75cm]
\]

\[
    \textbf{Όμως, έχουμε: } \hspace{0.25cm} 
    \lim_{x \to \infty} \left( \frac{\cos(x)}{x} \right) 
    \hspace{0.15cm} = \hspace{0.15cm} 0 \hspace{0.375cm} \textbf{από το ερώτημα 2!}
    \hspace{5.55cm} { }
    \\[0.75cm]
\]

\[
    \textbf{Και: } \hspace{0.25cm} 
    \lim_{x \to \infty} \left( \frac{3 \sqrt{x}}{x} \right) 
    \hspace{0.125cm} = \hspace{0.125cm}
    \lim_{x \to \infty} \left( \frac{3 \textbf{ · } \sqrt{x}}{\sqrt{x} \textbf{ · } \sqrt{x}} \right)
    \hspace{0.125cm} = \hspace{0.125cm}
    \lim_{x \to \infty} \left( \frac{3}{\sqrt{x}} \right)
    \hspace{0.125cm} = \hspace{0.125cm}
    \frac{3}{\infty} \hspace{0.125cm} = \hspace{0.125cm} 0
    \hspace{4.5cm} { }
    \\[0.75cm]
\]

\[
    \textbf{Άρα είναι: } \hspace{0.25cm} 
    \lim_{x \to \infty} \left( \frac{1 + 0 + 0}{1 + 0} \right) 
    \hspace{0.125cm} = \hspace{0.125cm}
    \frac{1}{1} \hspace{0.125cm} = \hspace{0.125cm} \boxed{1}
    \hspace{8.75cm} { }
    \\[10cm]
\]

% ---------------------------------------------------------------------------------------

\newpage

\vspace{5cm}

\begin{center}
    \noindent \boxed{\textbf{· 2\textsuperscript{η} Άσκηση:} \hspace{0.25cm} \textlatin{L} \hspace{0.10cm} = \hspace{0.10cm} \lim_{x \to x_0} \textlatin{f(x)} \hspace{0.125cm} \textbf{και} \hspace{0.125cm} 0 < |x - x_0| < \textbf{δ} \hspace{0.125cm} \Rightarrow \hspace{0.125cm} |\textlatin{}{f(x) - L}| < \hspace{0.125cm} \textbf{ε} \hspace{0.125cm} !}
\end{center}

\vspace{1.5cm}

\noindent \textbf{· Άσκηση 2\textsuperscript{η} - Ερώτημα 1\textsuperscript{ο}:} \hfill

\vspace{1.01cm}

\[
    \textbf{Έχουμε:} \hspace{0.25cm} f(x) = \frac{x^2 - 4}{x - 2}, \quad x_0 = 2, \quad \epsilon = 0.05.
    \hspace{7.99cm} { }
\]

\[
    \textbf{άρα:} \hspace{0.25cm} L \hspace{0.125cm} = \hspace{0.125cm} 
    \lim_{x \to 2} \left( \frac{x^2 - 4}{x - 2} \right)
    \hspace{0.125cm} = \hspace{0.125cm} \lim_{x \to 2} \left( \frac{x^2 - 2^2}{x - 2} \right)
    \hspace{0.125cm} = \hspace{0.125cm} \lim_{x \to 2} \left( \frac{(x - 2) \textbf{ · } (x + 2)}{x - 2} \right)
    \hspace{0.125cm} = \hspace{0.125cm} \lim_{x \to 2} \left( x + 2 \right)
    \hspace{0.125cm} = \hspace{0.125cm} 4 \hspace{0.1cm} { }
\]

\[
    \textbf{και:} \hspace{0.25cm} 0 < |x - 2| < \textbf{δ} \hspace{0.125cm} = 
    \hspace{0.125cm} |\frac{x^2 - 4}{x - 2} - 4| < \hspace{0.125cm} 0.05 \hspace{0.125cm} \Rightarrow
    \hspace{0.125cm} 0 < |x - 2| < \textbf{δ} \hspace{0.125cm} =
    \hspace{0.125cm} |x + 2 - 4| < \hspace{0.125cm} 0.05 \hspace{0.125cm} \Rightarrow
    \hspace{0.8cm} { }
\]

\[
    \Rightarrow \hspace{0.125cm} |x - 2| < \textbf{δ} \hspace{0.125cm} = 
    \hspace{0.125cm} |x - 2| < \hspace{0.125cm} 0.05 \hspace{0.35cm}
    \textbf{άρα:} \hspace{0.25cm} \boxed{\textbf{δ} \hspace{0.125cm} = \hspace{0.125cm} 0.05}
    \hspace{0.25cm} \textbf{ή} \hspace{0.25cm} \textbf{δ} \hspace{0.125cm} 
    = \hspace{0.125cm} \textbf{ε} \hspace{0.125cm} = \hspace{0.125cm} 0.05
    \hspace{3.75cm} { }
\]

% ---------------------------------------------------------------------------------------

\vspace{2.5cm}

\noindent \textbf{· Άσκηση 2\textsuperscript{η} - Ερώτημα 2\textsuperscript{ο}:} \hfill

\vspace{1.01cm}

\[
    \textbf{Έχουμε:} \hspace{0.25cm} f(x) = \sqrt{8 - 2x}, \quad x_0 = 2, \quad \epsilon = 0.1 
    \hspace{0.25cm} (\textlatin{Af}=(-\infty, 4]).
    \hspace{4.99cm} { }
\]

\[
    \textbf{άρα:} \hspace{0.25cm} L \hspace{0.125cm} = \hspace{0.125cm} 
    \lim_{x \to 2} \left( \sqrt{8 - 2x} \right)
    \hspace{0.125cm} = \hspace{0.125cm} \sqrt{4}
    \hspace{0.125cm} = \hspace{0.125cm} 2
    \hspace{8.5cm} { }
\]

\[
    \textbf{και έχουμε:} \hspace{0.25cm} 0 < |x - 2| < \textbf{δ} \hspace{0.125cm} \Rightarrow
    \hspace{0.125cm} |x - 2| < \textbf{δ} \hspace{0.125cm} \Rightarrow \hspace{0.125cm}
    - \textbf{δ} < x - 2 < \textbf{δ} \hspace{0.125cm} \Rightarrow \hspace{0.125cm}
    \hspace{3.45cm} { }
\]

\[
    \Rightarrow \hspace{0.125cm}
    2 - \textbf{δ} < x - 2 + 2 < 2 + \textbf{δ} \hspace{0.125cm} \Rightarrow \hspace{0.125cm}
    \Rightarrow \hspace{0.125cm} 2 - \textbf{δ} < x < 2 + \textbf{δ} \hspace{0.25cm} (1)
    \hspace{0.65cm} { }
\]

\[
    \textbf{και όμως:} \hspace{0.25cm} |f(x) - L| < \textbf{ε} \hspace{0.125cm} \Rightarrow \hspace{0.125cm}
    |\sqrt{8 - 2x} - 2| < 0.1 \hspace{0.125cm} \Rightarrow \hspace{0.125cm}
    -0.1 < \sqrt{8 - 2x} - 2 < 0.1 \hspace{0.125cm} \Rightarrow \hspace{0.125cm}
    \hspace{0.875cm} { }
\]

\[
    \Rightarrow \hspace{0.125cm}
    1.99 < \sqrt{8 - 2x} < 2.1 \hspace{0.125cm} \overset{(\geq 0 \forall x)}{\Rightarrow} \hspace{0.125cm}
    1.99^2 < 8 - 2x < 2.1^2 \hspace{0.125cm} \Rightarrow \hspace{0.125cm}
    1.99^2 - 8 < - 2x < 2.1^2 - 8 \hspace{0.125cm} \Rightarrow \hspace{0.125cm}
    \hspace{0.65cm} { }
\]

\[
    \overset{\div(-2)}{\Rightarrow} \hspace{0.125cm} 
    \frac{1.99^2 - 8}{-2} > x > \frac{2.1^2 - 8}{-2} \hspace{0.25cm} 
    \textbf{ή} \hspace{0.25cm} \frac{2.1^2 - 8}{-2} < x < \frac{1.99^2 - 8}{-2} \hspace{0.25cm} (2)
    \\[0.15cm]
\]

\[
    \textbf{άρα, από (1) και (2) έχουμε:} \\[0.05cm]
    \hspace{8.95cm} { }
\]

\[
    \textlatin{i.} \hspace{0.99cm}
    2 - \textbf{δ} \hspace{0.125cm} = \hspace{0.125cm} \frac{2.1^2 - 8}{-2}
    \hspace{0.125cm} \Rightarrow \hspace{0.125cm} \textbf{δ} 
    \hspace{0.125cm} \approx \hspace{0.125cm} 0,205
\]

\[
    \textlatin{ii.} \hspace{0.45cm}
    2 + \textbf{δ} \hspace{0.125cm} = \hspace{0.125cm} \frac{1.99^2 - 8}{-2}  
    \hspace{0.125cm} \Rightarrow \hspace{0.125cm} \textbf{δ} 
    \hspace{0.125cm} \approx \hspace{0.125cm} 0,01995
    \\[0.25cm]
\]

\[
    \textbf{έτσι, κρατάμε το:} \hspace{0.35cm} \boxed{ \textbf{δ} \hspace{0.125cm} \approx \hspace{0.125cm} 0,01995}
    \hspace{8.35cm} { }
\]

% ---------------------------------------------------------------------------------------

\newpage

\vspace{2.5cm}

\noindent \textbf{· Άσκηση 2\textsuperscript{η} - Ερώτημα 3\textsuperscript{ο}:} \hfill

\vspace{1.1cm}

\[
    \textbf{Έχουμε:} \hspace{0.25cm} f(x) = \frac{1}{x^2}, \quad x_0 = \sqrt{3}, \quad \epsilon = c > 0.
    \hspace{11.35cm} { } \\[0.25cm]
\]
    
\[
    \textbf{άρα:} \hspace{0.25cm} L \hspace{0.125cm} = \hspace{0.125cm} 
    \lim_{x \to \sqrt{3}} \left( \frac{1}{x^2} \right)
    \hspace{0.125cm} = \hspace{0.125cm} \frac{1}{(\sqrt{3})^2}
    \hspace{0.125cm} = \hspace{0.125cm} \frac{1}{3}
    \hspace{12.05cm} { } \\[0.25cm]
\]

\[
    \textbf{και έχουμε:} \hspace{0.35cm} 0 < |x - \sqrt{3}| < \textbf{δ} \hspace{0.125cm} \Rightarrow
    \hspace{0.125cm} \sqrt{3} - \textbf{δ}  < \hspace{0.125cm} x \hspace{0.125cm} < \sqrt{3} + \textbf{δ} 
    \hspace{0.35cm} (1)
    \hspace{8.15cm} { } \\[0.25cm]
\]

\[
    \textbf{και όμως:} \hspace{0.25cm} |f(x) - L| < \textbf{ε} \hspace{0.125cm} \Rightarrow \hspace{0.125cm}
    |\frac{1}{x^2} - \frac{1}{3}| < c \hspace{0.125cm} \Rightarrow \hspace{0.125cm}
    -c < \frac{1}{x^2} - \frac{1}{3} < c \hspace{0.125cm} \Rightarrow \hspace{0.125cm}
    \frac{1}{3} - c < \frac{1}{x^2} < \frac{1}{3} + c \hspace{0.125cm} \Rightarrow
    \hspace{3.55cm} { } \\[0.25cm]
\]

\[
    \overset{(-1)}{\underset{(c \approx 0)}{\Rightarrow}} \hspace{0.125cm}
    \frac{1}{\frac{1}{3} - c} > x^2 > \frac{1}{\frac{1}{3} + c}
    \hspace{0.125cm} \overset{(\sqrt{})}{\Rightarrow} \hspace{0.125cm}
    \sqrt{\frac{1}{\frac{1}{3} - c}} > |x| > \sqrt{\frac{1}{\frac{1}{3} + c}}
    \hspace{0.125cm} \Rightarrow \hspace{0.125cm} 
    \sqrt{\frac{1}{\frac{1}{3} + c}} < |x| < \sqrt{\frac{1}{\frac{1}{3} - c}}
    \hspace{0.125cm} \Rightarrow \hspace{0.125cm} 
    \hspace{3.55cm} { } \\[0.25cm]
\]

\[
    \textbf{αλλά:} \hspace{0.35cm} x \neq 0 \hspace{0.35cm} \textbf{οπότε έχουμε τις περιπτώσεις:} \hspace{0.35cm}
    \hspace{9.65cm} { } \\[0.25cm]
\]

\[
    \textlatin{i: } \hspace{0.35cm} x > 0 \hspace{0.35cm} \textbf{:} \hspace{0.35cm}
    \sqrt{\frac{1}{\frac{1}{3} + c}} < x < \sqrt{\frac{1}{\frac{1}{3} - c}} \hspace{0.35cm} (2)
    \hspace{11.45cm} { } \\[0.25cm]
\]

\[
    \textbf{και από (1) - (2):} \hspace{0.35cm} \sqrt{3} \hspace{0.125cm} - \hspace{0.125cm} \textbf{δ} \hspace{0.125cm} = \hspace{0.125cm} \sqrt{\frac{1}{\frac{1}{3} + c}} \hspace{0.135cm} \Rightarrow \hspace{0.135cm} 
    \sqrt{3} \hspace{0.125cm} = \hspace{0.125cm} \sqrt{\frac{1}{\frac{1}{3} + c}} \hspace{0.125cm} + \hspace{0.125cm} \textbf{δ} \hspace{0.135cm} \overset{(>0)}{\Rightarrow} \hspace{0.135cm} 
    3 \hspace{0.125cm} = \hspace{0.125cm} \frac{1}{\frac{1}{3} + c} \hspace{0.125cm} + \hspace{0.125cm} \delta^2
    \hspace{0.125cm} \Rightarrow \hspace{0.125cm} \delta = \sqrt{3 - \frac{1}{\frac{1}{3} + c}}  \\[0.25cm]
\]

\[
    \textlatin{ii: } \hspace{0.35cm} x < 0 \hspace{0.35cm} \textbf{:} \hspace{0.35cm}
    \sqrt{\frac{1}{\frac{1}{3} + c}} < -x < \sqrt{\frac{1}{\frac{1}{3} - c}} \hspace{0.35cm}
    \Rightarrow \hspace{0.35cm} - \sqrt{\frac{1}{\frac{1}{3} + c}} > x > - \sqrt{\frac{1}{\frac{1}{3} - c}} 
    \hspace{0.35cm} \Rightarrow \hspace{0.35cm}
    - \sqrt{\frac{1}{\frac{1}{3} - c}} < x < - \sqrt{\frac{1}{\frac{1}{3} + c}} \hspace{0.35cm} (2)  \\[0.25cm]
\]

\[
    \textbf{και από (1) - (2):} \hspace{0.35cm} \sqrt{3} + \delta \hspace{0.125cm} = \hspace{0.125cm}
    - \sqrt{\frac{1}{\frac{1}{3} + c}} \hspace{0.35cm} \Rightarrow \hspace{0.35cm}
    \delta = -(\sqrt{\frac{1}{\frac{1}{3} + c}} + \sqrt{3})
    \hspace{6.45cm} { } \\[0.25cm] \\[0.25cm]
\]

\[
    \textbf{έτσι, καταλήγουμε:} \hspace{0.35cm} \delta \hspace{0.125cm} = \hspace{0.125cm}
    min(-(\sqrt{\frac{1}{\frac{1}{3} + c}} + \sqrt{3}), \quad \sqrt{3 - \frac{1}{\frac{1}{3} + c}}) 
    \hspace{0.35cm} \textbf{ανάλογα το \textlatin{c} ...}
    \hspace{3.75cm} { } \\[0.25cm]
\]

% ---------------------------------------------------------------------------------------

\newpage

\vspace{5cm}

\begin{center}
    \noindent \boxed{\textbf{· 3\textsuperscript{η} Άσκηση:} \hspace{0.25cm} \textbf{Να βρεθούν οι ασύµπτωτες}!}
\end{center}

\vspace{1.5cm}

\noindent \textbf{· Άσκηση 3\textsuperscript{η} - Ερώτημα 1\textsuperscript{ο}:} \hfill

\vspace{0.75cm}

\[
    \textbf{Έχουμε:} \hspace{0.35cm} y = \frac{x^2 + 2x - 3}{x^2 - x - 6} \hspace{0.35cm} \textbf{ή} \hspace{0.35cm}
    \textlatin{f(x)} = \frac{x^2 + 2x - 3}{x^2 - x - 6} \hspace{0.35cm} \textbf{με} \hspace{0.35cm} A_f = {R} - \{3, -2\} .
    \hspace{0cm} { } \hfill
\]

\[
    \textbf{Για κατακόρυφες ασύμπτωτες:} \hspace{6.35cm} { } 
\]

\[
\lim_{x \to -2} f(x) = \lim_{x \to -2} \frac{(-2)^2 + 2(-2) - 3}{(-2)^2 - (-2) - 6} = \infty
\]

\[
    \textbf{και:}
\]

\[
\lim_{x \to 3} f(x) = \lim_{x \to 3} \frac{(3)^2 + 2(3) - 3}{(3)^2 - (3) - 6} = \infty
\]

\[
    \textbf{Άρα οι κατακόρυφες ασύμπτωτες είναι: } \quad \boxed{x = -2} \quad \textbf{και} \quad \boxed{x = 3}
\]

\[
    \textbf{Για οριζόντιες ασύμπτωτες:} \hspace{6.35cm} { }
\]

\[
    \lim_{x \to \infty} f(x) = \lim_{x \to \infty} \frac{x^2 + 2x - 3}{x^2 - x - 6} = \lim_{x \to \infty} \frac{x^2}{x^2} = \lim_{x \to \infty} 1 = 1
\]

\[
    \textbf{Άρα η οριζόντια ασύμπτωτη είναι: } \quad \boxed{y = 1} \hspace{3.75cm} { }
\]

\[
    \textbf{Και αφού υπάρχει οριζόντια ασύμπτωτη δεν υπάρχει πλάγια!}
\]


% -----------------

\vspace{2.25cm}

\noindent \textbf{· Άσκηση 3\textsuperscript{η} - Ερώτημα 2\textsuperscript{ο}:} \hfill

\vspace{0.75cm}

\[
    \textbf{Έχουμε:} \hspace{0.35cm} y = \sqrt{\frac{x^2 + 4}{4x^2 + 1}} \hspace{0.35cm} \textbf{ή} \hspace{0.35cm}
    \textlatin{f(x)} = \sqrt{\frac{x^2 + 4}{4x^2 + 1}} \hspace{0.35cm} \textbf{με} \hspace{0.35cm} A_f = {R}.
    \hspace{2.55cm} { }
\]

\[
    \textbf{Αφού} \quad A_f = R \quad \textbf{δεν υπάρχουν κατακόρυφες ασύμπτωτες}. \hspace{1.85cm} { }
\]
\[
    \textbf{Για οριζόντιες ασύμπτωτες έχουμε:} \hspace{5.75cm} { }
\]

\[
    \lim_{x \to \infty} f(x) = \lim_{x \to \infty} \sqrt{\frac{x^2 + 4}{4x^2 + 1}} = \sqrt{\lim_{x \to \infty} \frac{x^2}{4x^2}} = \sqrt{\lim_{x \to \infty} \frac{1}{4}} = \sqrt{ \frac{1}{4}} = \frac{\sqrt{1}}{\sqrt{4}} = \frac{1}{2}
\]

\[
    \textbf{Άρα η οριζόντια ασύμπτωτη είναι: } \quad \boxed{y = \frac{1}{2}} \hspace{4.1cm} { }
\]

\[
    \textbf{Και αφού υπάρχει οριζόντια ασύμπτωτη δεν θα υπάρχει πλάγια!}
\]


% -----------------

\newpage

\vspace{5cm}

\noindent \textbf{· Άσκηση 3\textsuperscript{η} - Ερώτημα 3\textsuperscript{ο}:} \hfill

\vspace{0.75cm}


\[
    \textbf{Έχουμε:} \hspace{0.35cm} y = \frac{3x^\frac{3}{2} + x}{\sqrt{x} - 2} \hspace{0.35cm} 
    \textbf{ή} \hspace{0.35cm}
    \textlatin{f(x)} = \frac{3x^\frac{3}{2} + x}{\sqrt{x} - 2} \hspace{0.35cm} \textbf{με} \hspace{0.35cm} 
    A_f = [0, 4) \cup (4, +\infty). \hspace{5.5cm} { }
\]
    
\[
    \textbf{Για κατακόρυφες ασύμπτωτες:} \hspace{11.25cm} { }
\]

\[
    \textlatin{i:} \quad
    \lim_{x \to 0} f(x) = \lim_{x \to 0} \frac{3x^\frac{3}{2} + x}{\sqrt{x} - 2} =
    \frac{0}{- 2} = 0 \\[0.25cm] \hspace{4.05cm} { }
\]

\[
    \textbf{Αφού} \quad \lim_{x \to 0} f(x) \neq \infty \quad \textbf{δεν υπάρχει κατακόρυφη ασύμπτωτη στο} x=0. \hspace{4.35cm} { }
\]

\[
    \textlatin{ii:} \quad
    \lim_{x \to 4^-} f(x) = \lim_{x \to 4^-} \frac{3x^\frac{3}{2} + x}{\sqrt{x} - 2} =
    \lim_{x \to 4^-} \frac{3(\sqrt{4^3}) + 4}{\sqrt{4} - 2} =
    \frac{28}{0^-} = - \infty
\]

\[  
    \textbf{και:} \quad
    \lim_{x \to 4^+} f(x) = \lim_{x \to 4^+} \frac{3x^\frac{3}{2} + x}{\sqrt{x} - 2} =
    \lim_{x \to 4^+} \frac{3(\sqrt{4^3}) + 4}{\sqrt{4} - 2} =
    \frac{28}{0^+} = + \infty \\[0.25cm] \hspace{0.35cm} { }
\]

\[
    \textbf{Αφού} \quad \lim_{x \to 4} f(x) = \infty \quad \textbf{έχουμε κατακόρυφη ασύμπτωτη την} \quad \boxed{x=4}. \hspace{4.35cm} { }
\]

\[
    \textbf{Για οριζόντιες ασύμπτωτες έχουμε:} \hspace{9.85cm} { }
\]

\[
    \lim_{x \to +\infty} f(x) = \lim_{x \to +\infty} \frac{3x^\frac{3}{2} + x}{\sqrt{x} - 2} = 
    \lim_{x \to +\infty} \frac{3x^\frac{3}{2} + x^1}{x^\frac{1}{2} - 2} = 
    \lim_{x \to +\infty} \frac{3x^\frac{3}{2}}{x^\frac{1}{2}} = 
    \lim_{x \to +\infty} \frac{3x^\frac{3}{2}}{x^\frac{1}{2}} = 
    \lim_{x \to +\infty} 3x^\frac{3}{2} \textbf{ · } x^{-\frac{1}{2}} = 
    \lim_{x \to +\infty} 3x = +\infty \\[0.25cm]
\]

\[
    \textbf{Αφού} \quad \lim_{x \to \infty} f(x) \neq l \hspace{0.125cm} \textbf{ε} \hspace{0.125cm} R \quad \textbf{δεν υπάρχει οριζόντια ασύμπτωτη!} \hspace{5.15cm} { }
\]

\[
    \textbf{Άρα, για πλάγιες ασύμπτωτες θα είναι:} \hspace{8.7cm} { }
\]

\[
    \textlatin{i:} \quad
    \lim_{x \to \infty} \frac{f(x)}{x} = \lambda \implies 
    \lim_{x \to \infty} \left( \frac{\frac{3x^{\frac{3}{2}} + x}{\sqrt{x} - 2}}{\frac{x}{1}}
    \right) = \lim_{x \to \infty} \left( \frac{3x^{\frac{3}{2}} + x}{x(\sqrt{x} - 2)}  
    \right) = \lim_{x \to \infty} \left( \frac{3x^{\frac{3}{2}} + x^1}{x^1 \textbf{ · } x^{\frac{1}{2}} - 2x} \right) \quad \Rightarrow
\]  
    
\[
    \overset{\textbf{Ιδ.}}{\Rightarrow} \quad \lim_{x \to \infty} \left( \frac{3x^{\frac{3}{2}} + x^1}{x^{\frac{2}{2} + \frac{1}{2}} - 2x^1} \right) = \lim_{x \to \infty} \left( \frac{3x^{\frac{3}{2}} + x^1}{x^{\frac{3}{2}} - 2x^1} \right)
    \overset{\textbf{Ιδ.}}{=} \lim_{x \to \infty} \left( \frac{3x^{\frac{3}{2}}}{x^{\frac{3}{2}}} \right)
    = \lim_{x \to \infty} \left( 3 \right) = 3 \quad \Rightarrow \quad \underline{\lambda = 3}
\]

\[
    \textlatin{ii:} \quad
    \lim_{x \to \infty} f(x) - \lambda x = \beta \implies 
    \lim_{x \to \infty} \left( \frac{3x^{\frac{3}{2}} + x}{\sqrt{x} - 2} - 3x
    \right) = \lim_{x \to \infty} \left( \frac{3x^{\frac{3}{2}} + x}{\sqrt{x} - 2} 
    - \frac{3x(\sqrt{x} -2)}{\sqrt{x} -2}  \right) 
    \quad \Rightarrow
\]  

\[
    \Rightarrow
    \lim_{x \to \infty} \left( \frac{3x^{\frac{3}{2}} + x}{\sqrt{x} - 2} 
    - \frac{3x(\sqrt{x} -2)}{\sqrt{x} -2}  \right)
    = \lim_{x \to \infty} \left( \frac{3x^{\frac{3}{2}} + x}{\sqrt{x} - 2} 
    - \frac{3x^{\frac{3}{2}} - 6x)}{\sqrt{x} -2}  \right)
    = \lim_{x \to \infty} \left( \frac{7x}{\sqrt{x} - 2}\right)
    = \lim_{x \to \infty} \left( \frac{7x}{x^{\frac{1}{2}}}\right)
    = \lim_{x \to \infty} ( 7 \sqrt{x} ) = \infty \\[0.25cm]
\]

\[
    \textbf{Άρα} \quad \underline{\beta = +\infty} \quad \textbf{αδύνατο, οπότε δεν υπάρχουν ούτε πλάγιες ασύμπτωτες!}  \hspace{2.55cm} { }
\]

% ---------------------------------------------------------------------------------------

\newpage

\vspace{10cm}

\begin{center}
    \noindent \boxed{\textbf{· 4\textsuperscript{η} Άσκηση:} \hspace{0.25cm} \textbf{Θεώρημα} \hspace{0.175cm} \textlatin{Bolzano} \quad!}
\end{center}

\vspace{1.5cm}

\noindent \textbf{· Η λύση μου:} \hfill

\vspace{2.25cm}

\[
    \textbf{Έχουμε την εξίσωση:} \quad \sin(x) = x - 1 \quad \textbf{ή αλλιώς:} \quad 
    \textbf{ημ(\textlatin{x})} = x - 1 \quad \hspace{4.99cm} { }
\]

\[
    \textbf{Θέτω:} \quad f(x) = \sin(x) - x + 1 
    \quad \textbf{με} \quad \textlatin{Af} = \mathbb{R} \hspace{9.55cm} { }
\]

\[
    \textbf{Και:} \quad \textlatin{i.} \quad f(0) = \sin(0) - 0 + 1 
    \hspace{0.125cm} = \hspace{0.125cm} 0 - 0 + 1 \hspace{0.125cm} = \hspace{0.125cm} 1 \hspace{0.125cm} > \hspace{0.125cm} 0 \hspace{7.75cm} { }
\]

\[
    \textlatin{ii.} \quad f(1) = \sin(1) - 1 + 1 
    \hspace{0.125cm} \approx \hspace{0.125cm} 0.84147 - 1 + 1 \hspace{0.125cm} \approx \hspace{0.125cm} 0.84147 \hspace{0.125cm} > \hspace{0.125cm} 0 \hspace{4.45cm} { }
\]

\[
    \textlatin{iii.} \quad f(2) = \sin(2) - 2 + 1 
    \hspace{0.125cm} \approx \hspace{0.125cm} 0.90929 - 2 + 1 \hspace{0.125cm} \approx \hspace{0.125cm} - 0.09071 \hspace{0.125cm} < \hspace{0.125cm} 0 \hspace{4.05cm} { }
\]

\[
    \textbf{Εφαρμόζω θεώρημα \textlatin{Bolzano} στην \textlatin{f}(\textlatin{x}) στο διάστημα } 
    \hspace{0.1cm} [0, 2]: \hspace{5.25cm} { }
\]

\[
    \implies f \textbf{ συνεχής στο } [0, 2] \textbf{ ως πράξεις συνεχών συναρτήσεω.}
\]

\[
    \implies f(0) > 0 \textbf{ και } f(2) < 0 \textbf{   άρα   } f(0) \textbf{ · } f(2) < 0 
    \hspace{3.3cm} { }
\]

\[
    \textbf{Έτσι, υπάρχει τουλ. ένα } x_0 \textbf{ ε } (0, 2) \textbf{ τέτοιο ώστε } f(x_0) = 0
    \hspace{6.15cm} { }
\]

\[
    \textbf{και άρα η εξίσωση } \sin(x) = x - 1 \textbf{ να έχει τουλ. μία λύση σε ένα συγκεκριμένο διάστημα!}
\]


% ---------------------------------------------------------------------------------------

\newpage

\vspace{10cm}

\begin{center}
    \noindent \boxed{\textbf{· 5\textsuperscript{η} Άσκηση:} \hspace{0.25cm} \textbf{Εφαπτώμένες και Κάθετες Ευθείες!}}
\end{center}

\vspace{1.5cm}

\noindent \textbf{· Άσκηση 5\textsuperscript{η} - Ερώτημα 1\textsuperscript{ο}:} \hfill

\vspace{0.85cm}

\[
    \textbf{Έχουμε: } \quad x^2 + xy - y^2 = 1 \quad \textbf{ και } \quad \textbf{ Δ} = (2, 3) = (x_0, y_0) \quad \hspace{5.01cm} { }
    \\[0.35cm]
\]

\[
    \textbf{Άρα θα είναι: } \quad (x^2)' + (xy)' - (y^2)' = (1)' \quad \textbf{ ή } \quad 
    \frac{d}{dx}(x^2) + \frac{d}{dx}(xy) - \frac{d}{dx}(y^2) = \frac{d}{dx}(1) \quad \Rightarrow \quad
\]

\[ 
    \hspace{0.75cm} { }
    \implies \quad 2x + y + xy' - 2y'y = 0 \quad \textbf{ ή } \quad 
    2x + y - x\frac{dy}{dx} - 2y\frac{dy}{dx} = 0 \quad \Rightarrow \quad
\]

\[ 
    \hspace{1.45cm} { }
    \implies \quad 2x + y = 2y'y - xy' \quad \textbf{ ή } \quad 
    2x + y = 2y\frac{dy}{dx} - x\frac{dy}{dx} = 0 \quad \Rightarrow \quad
\]

\[ 
    \hspace{0.85cm} { }
    \implies \quad 2x + y = 2y'y - xy' \quad \textbf{ ή } \quad 
    2x + y = 2y\frac{dy}{dx} - x\frac{dy}{dx} \quad \Rightarrow \quad
\]

\[ 
    \hspace{0.85cm} { }
    \implies \quad 2y'y - xy' = 2x + y \quad \textbf{ ή } \quad 
    2y\frac{dy}{dx} - x\frac{dy}{dx} = 2x + y \quad \Rightarrow \quad
\]

\[ 
    \hspace{0.85cm} { }
    \implies \quad y' (2y - x) = 2x + y \quad \textbf{ ή } \quad 
    \frac{dy}{dx} ( 2y - x) = 2x + y \quad \Rightarrow \quad
\]

\[ 
    \hspace{0.85cm} { }
    \implies \quad y' = \frac{2x + y}{2y - x} \quad \textbf{ ή } \quad 
    \frac{dy}{dx} = \frac{2x + y}{2y - x} \quad \Rightarrow \quad
\]

\[ 
    \implies \quad y' = \frac{2x + y}{2y - x} \quad \textbf{ ή } \quad 
    \frac{dy}{dx} = \frac{2x + y}{2y - x} \quad \\[0.15cm]
\]

\[ 
    \textbf{άρα: } \quad
    y' \big|_{(2,3)} = \frac{2(2) + 3}{2(3) - 2} = \frac{7}{4} \quad \textbf{ ή } \quad 
    \frac{dy}{dx} \big|_{(2,3)} = \frac{2(2) + 3}{2(3) - 2} = \frac{7}{4} \\[0.15cm] \hspace{0.15cm} { }
\]

\[ 
    \textbf{οπότε } \quad
    \boxed{y' = \frac{7}{4}}
\]

\[ 
    \textbf{και η εξίσωση εφαπτωμένης είναι: \quad εφ:} \quad
    y - y_0 = y'(x - x_0) \hspace{3.35cm} { }
\]

\[ 
    \implies \quad y - 3 = \frac{7}{4} (x - 2) \quad \Rightarrow \quad \boxed{y = \frac{7}{4} x - \frac{1}{2}} \\[0.15cm]
\]

\[ 
    \textbf{και για την κάθετη ευθεία γνωρίζουμε την κλίση της:} \quad \lambda_\epsilon \perp \lambda_{\textbf{εφ}} \hspace{2.15cm} { }
\]

\[ 
    \implies \quad \lambda_\epsilon = - \frac{1}{ \lambda_{\textbf{εφ}} } \quad
    \Rightarrow \quad \lambda_\epsilon = - \frac{1}{ \frac{7}{4} } = \boxed{ - \frac{4}{7} = \lambda_\epsilon} \\[0.15cm]
\]

\[ 
    \textbf{και η κάθετη ευθεία θα είναι: \quad ε:} \quad
    y - y_0 = \lambda_\epsilon(x - x_0) \quad \Rightarrow \quad
    \boxed{y = - \frac{4}{7}x + \frac{29}{7}} \hspace{0.35cm} { }
\]

% ---------------------------------------------------------------------------------------

\newpage

\vspace{1.5cm}

\noindent \textbf{· Άσκηση 5\textsuperscript{η} - Ερώτημα 2\textsuperscript{ο}:} \hfill

\vspace{0.75cm}

\[
    \textbf{Έχουμε: } \quad x \sin(2y) = y \cos(2x) \quad \textbf{ και } \quad 
    \textbf{ Δ} = (\frac{\pi}{4}, \frac{\pi}{2}) = (x_0, y_0) \quad
    \\[0.35cm] \hspace{8.25cm} { }
\]

\[
    \textbf{Άρα: } \quad 
    \\ (x \sin(2y))' = (y \cos(2x))' \quad
    \textbf{όμως ξέρουμε:} \quad (fg)' = f'g + fg' \quad
    \textbf{και} \quad \sin(u) ' = \cos(u) u', \quad \cos(u)' = - \sin(u) u'
\]

\[
    \implies \quad \sin(2y) + x 2 \cos(2y) y' = \cos(2x) y' - 2y \sin(2x) \quad \Rightarrow \quad 
    2x \cos(2y)y' - \cos(2x)y' = - \sin(2y) - 2y \sin(2x) \quad \Rightarrow
\]

\[
    \implies \quad y' (2x\cos(2y) - \cos(2x) ) = - \sin(2y) - 2y \sin(2x) \sin\quad \Rightarrow \quad 
    y' = - \frac{\sin(2y) + 2y\sin(2x)}{2x\cos(2y) - \cos(2x)} \\[0.15cm]
\]

\[
    \textbf{Επίσης γνωρίζουμε: } \quad 
    \sin{(\pi)} = 0, \quad \sin{(\frac{\pi}{2})} = 1, \quad \cos{(\pi)} = -1, \quad \cos{(\frac{\pi}{2})} = 0.
    \\[0.15cm] \hspace{5.35cm} { }
\]

\[
    \textbf{άρα: } \quad 
    y' \big|_{(\frac{\pi}{4},\frac{\pi}{2})} = - \frac{\sin(\pi) + \pi \sin(\frac{\pi}{2})}
    {\frac{\pi}{2} \cos(\pi) - \cos(\frac{\pi}{2})} \quad = \quad
    \frac{-\pi}{-\frac{\pi}{2}} \quad = \quad \underline{2} \\[0.25cm] \hspace{9.15cm} { }
\]

\[
    \textbf{οπότε είναι: } \quad 
    \textbf{εφ: } \quad y - y_0 = y'(x - x_0) \quad \Rightarrow \quad
    y - \frac{\pi}{2} = 2 (x - \frac{\pi}{4}) \quad \Rightarrow \quad
    \boxed{y = 2x} \\[0.15cm] \hspace{4.95cm} { }
\]

\[
    \textbf{και θα έχουμε: } \quad 
    \lambda_\epsilon = - \frac{1}{y'} = - \frac{1}{2} \quad \Rightarrow \quad \\[0.15cm] \hspace{11.65cm} { }
\]

\[
    \textbf{έτσι: } \quad 
    \textbf{ε: } \quad y - y_0 = \lambda_\epsilon(x - x_0) \quad \Rightarrow \quad
    y - \frac{\pi}{2} = - \frac{1}{2} (x - \frac{\pi}{4}) \quad \Rightarrow \quad
    \boxed{y = - \frac{1}{2} x + \frac{5 \pi}{8}} \\[0.25cm] \hspace{4.75cm} { }
\]

% ---------------------------------------------------------------------------------------

\vspace{1.5cm}

\noindent \textbf{· Άσκηση 5\textsuperscript{η} - Ερώτημα 3\textsuperscript{ο}:} \hfill

\vspace{0.75cm}

\[
    \textbf{Έχουμε: } \quad x^2 \cos^2(y) - \sin(y) = 0 \quad \textbf{ και } \quad 
    \textbf{ Δ} = (0, \pi) = (x_0, y_0) \quad  \textbf{αντίστοιχα.}
    \\[0.35cm] \hspace{3.4cm} { }
\]

\[
    \textbf{Άρα: } \quad 
    (x^2 \cos^2(y))' - (\sin(y))' = (0)' \quad \Rightarrow \quad
    2x\cos^2(y) - 2x^2\cos(y)\sin(y) y' - \cos(y) y' = 0 \quad \Rightarrow \quad \hspace{0.85cm} { }
\]

\[
    \implies \quad 2x\cos^2(y) = 2x^2\cos(y)\sin(y)y' + \cos(y)y' \quad \Rightarrow \quad 
    2x\cos^2(y) = y' (2x^2\cos(y)\sin(y) + \cos(y)) \quad \Rightarrow \quad 
\]

\[
    \implies \quad y' = \frac{2x\cos^2(y)}{(2x^2\cos(y)\sin(y) + \cos(y))} \quad \Rightarrow \quad 
    y' \big|_{(0, \pi)} = \frac{(-2)(0)(-1)^2}{(-2)(0^2)(0) + (-1)} \quad = \quad \boxed{0} \\[0.25cm]
     \hspace{0.25cm} { }
\]

\[
    \textbf{και: } \quad 
    \textbf{εφ: } \quad y - y_0 = y'(x - x_0) \quad \Rightarrow \quad
    y - \pi = 0 (x - 0) \quad \Rightarrow \quad
    \boxed{y = \pi} \\[0.15cm] \hspace{4.65cm} { }
\]

\[
    \underline{\textbf{Αλλά, δεν ορίζεται κάθετη ευθεία (ε)}} \textbf{ αφού: } \quad 
    \textbf{εφ: } \quad \lambda_\epsilon = - \frac{1}{y'} = - \frac{1}{0} \quad \hspace{3.45cm} { }
\]

% ---------------------------------------------------------------------------------------

\newpage

\begin{center}
    \vspace*{\fill}
    \textbf{\LARGE{ΤΕΛΟΣ}} \\[0.001cm]
    \textbf{(ευχαριστώ)}
    \\[5cm]
    \vspace*{\fill}
\end{center}

% ---------------------------------------------------------------------------------------

\end{document}

% ---------------------------------------------------------------------------------------
